\section{Discussion}\label{sec:discussion}

The two doors are not two labels for the same debt.
A compliance study that pools enrolled \vat{} certificates will let the assessed tail dominate any value-weighted moment and the declared mass dominate any count-weighted moment.
The literature on reporting and third-party information remains the right place to study how a balance is established \citep{kleven2011,pomeranz2015,naritomi2019}.
It is the wrong place to stop if the question is the stock the attorney office holds.
Declared debt is unpaid reported tax.
Assessed debt is tax the reporting system did not close, or a correction the administration imposed.
Once both are certificates, they still do not look alike.
Size, age and the relation between payment and the remaining balance all differ, and the differences are large enough that the pooled object is mostly the assessed object whenever the researcher weights by value.

Payment and recovery answer different administrative questions, which is why they should be reported together.
A rule that counts certificates with any payment will find the assessed segment ahead.
A rule that asks how many reais have been collected per real still enrolled will find the declared segment ahead.
The second result is what one would expect when assessed balances are much larger and, often, older and more likely to sit with a closed or unfit firm.
It is not evidence that declared debt is easy to collect in absolute reais, and it is not evidence that assessed debt is untouched.
Some payment on a large certificate can be small relative to the balance and still material in currency.
Collection statistics that use only one of the two rates will rank the doors in the order that rate produces.

The 1~percent list is a description of where the stock sits, not a treatment effect.
The Lorenz curve says that a specialised review of the largest certificates is a review of most of the money, and that the list is an assessed list with a declared remainder.
That is useful for the attorney office, because the alternative is a book of more than two million declared certificates whose combined value is the smaller stock.
It does not say that working the list raises recovery.
We do not observe the counterfactual, and we do not observe which certificates were already under active collection.
The within-segment curves are the warning against reading the pooled 75~percent as a general law.
Inside each door, the largest 1~percent hold about two fifths of that door's stock.
A manager who applies the pooled cutoff inside declared debt alone will not capture three quarters of declared value.

The firm patterns line up with the two doors without identifying a cause.
Assessed stock sits with older firms, with more closed and unfit registrations, and with a heavier weight on wholesale trade, telecommunications and a few large manufacturing divisions.
Declared stock sits with a wider set of firms, with more of the value in manufacturing, land transport and petroleum products, and with almost all of the value in firms that hold many certificates.
A natural reading is that assessment is aimed at larger, established operations, while declared debt is the unpaid remainder of ordinary filing, including many small certificates on the same firm.
The register is later than the debt, so closure and unfitness can be consequences.
We treat that reading as a description, not as a result this cross-section can establish.

For an administration outside Brazil, the useful object is the distinction rather than the S\~ao Paulo percentages.
A VAT system that can flag tax the taxpayer declared and did not remit, separately from tax fixed by an assessment, can report the same four moments: size, concentration, payment and recovery.
The European invoice-credit VAT is the framing for that distinction, not a claim that EU debt stocks have this shape \citep{keen2006}.

The linearity diagnostic selects different payment equations in the two segments, and it sets ordinary least squares aside for the balance.
The fitted associations are partial associations in that specification, on one cross-section.
They are not the effect of a collection rule.
\textbf{[DISCUTIR E ESCREVER: Redigir a leitura econômica dos coeficientes ajustados, inclusive a hipótese de que o declarado não consegue pagar e o autuado não quer. | Contexto: a discussão já estabeleceu que pagamento e recuperação ordenam as duas portas em sentidos diferentes e que as associações ajustadas não são o efeito de uma regra de cobrança; falta dizer se os coeficientes sustentam essa hipótese ou apenas descrevem o corte transversal. | Dados necessários: Table~\ref{tab:profile}, linhas de qualquer pagamento e de recuperação; Table~\ref{tab:logit-ame}, linha Assessed (declared reference), colunas AME (pp) e intervalos cluster e HC3; Figure~\ref{fig:origin-contrasts}, panel~(a); Table~\ref{tab:within-firm}, a mesma linha, colunas Difference (pp) e intervalos; Figure~\ref{fig:origin-contrasts}, panel~(b); Table~\ref{tab:recovery-pct}, a mesma linha, coluna Percent difference e intervalos; Figure~\ref{fig:origin-contrasts}, panel~(c); Table~\ref{tab:secsize-payment}, Table~\ref{tab:secsize-recovery}, Table~\ref{tab:secsize-quantile} e Table~\ref{tab:secsize-mixed}, linhas Manufacturing, Transportation and storage, Information and communication, Small e Other sizes. Crédito e liquidez da firma: not yet available. Histórico de protesto, além do status corrente: not yet available. Histórico completo de ajuizamento, além da cobertura parcial da tabela de ajuizamento: not yet available.]}
